Сообщения вкладки МКО выводятся в журнал операций. Журнал используется для контроля действий оператора и диагностики обмена.

В журнал выводятся:

\begin{itemize}
	\item факт готовности вкладки и адрес gRPC endpoint;
	\item параметры выполненных вызовов;
	\item сообщения об успешном выполнении команд;
	\item ошибки gRPC и ошибки открытия файлов;
	\item результаты чтения подадреса ОУ;
	\item события входящих команд ОУ из stream-подписки;
	\item сообщения о загрузке и сохранении СД из текстовых и двоичных файлов;
	\item предупреждения о лишних словах в файле или некратном 4 байтам размере двоичного файла;
	\item сообщения об автоматическом изменении номера МКО, если один номер был выбран одновременно для КК и ОУ.
\end{itemize}

\begin{table}[H]
	\caption{Сообщения вкладки МКО и действия оператора}
	\begin{tabularx}{\textwidth}{|p{0.40\textwidth}|X|}
		\hline
		\textbf{Сообщение или тип сообщения} & \textbf{Действие оператора} \\
		\hline
		Сообщение о готовности вкладки и адресе gRPC endpoint & Убедиться, что вкладка подключается к нужному микросервису. При неверном адресе проверить конфигурацию запуска. \\
		\hline
		Сообщение об успешном выполнении команды & Перейти к следующей операции по выбранному сценарию работы. \\
		\hline
		Ошибка gRPC & Проверить, запущен ли микросервис МКО, доступен ли драйвер МКО или заглушка драйвера, корректно ли указан адрес подключения. После устранения причины повторить вызов. \\
		\hline
		Ошибка проверки введенных значений & Проверить диапазоны значений, формат чисел, номер МКО, адрес ОУ, подадрес и количество слов данных. После исправления повторить команду. \\
		\hline
		Ошибка открытия или чтения файла & Проверить путь к файлу, права доступа, формат файла и его содержимое. После исправления повторить загрузку. \\
		\hline
		Предупреждение о лишних словах в файле & Учитывать, что загружены только первые 32 слова. При необходимости сократить файл и повторить загрузку. \\
		\hline
		Предупреждение о размере двоичного файла, не кратном 4 байтам & Учитывать, что неполный хвост файла был проигнорирован. При необходимости исправить файл и повторить загрузку. \\
		\hline
		Сообщение об автоматическом изменении номера МКО & Проверить выбранные номера МКО в режимах КК и ОУ. Один номер не должен одновременно использоваться в двух режимах. \\
		\hline
		Событие входящей команды ОУ & Проанализировать командное слово, направление обмена, адрес ОУ, подадрес, количество слов и расшифровку результата. \\
		\hline
		Результат чтения подадреса ОУ & Проверить подадрес, командное слово, слово результата, расшифровку результата и массив принятых слов данных. \\
		\hline
	\end{tabularx}
\end{table}

Если вызов завершился ошибкой, оператор должен проверить введенные значения, доступность микросервиса МКО, адрес платы и состояние заглушки драйвера или реального драйвера. После устранения причины сбоя допускается повторный запуск команды.
